\chapter{Conclusion and Future Work}
\label{ch:conclusion}

\section{Conclusion}
This thesis has presented the \projname: an end-to-end, effectively serverless
machine-learning system that forecasts the Air Quality Index up to 72 hours ahead
for 22 cities across every province and territory of Pakistan, retrains and
validates itself nightly without human intervention, explains each forecast in
plain language, quantifies its own uncertainty, monitors its own inputs and
outputs for decay, and serves all of this publicly at a running cost of
effectively zero.

The work began from the observation that air-quality information in Pakistan is
the wrong shape for the decisions people make with it. It is a nowcast when a
forecast is needed; it covers one or two cities when the worst-affected city in
the record assembled here is not the one that receives the attention; it is
offered without explanation or uncertainty; and where a model exists at all it is
a model trained once, which for a strongly seasonal phenomenon means a model that
decays silently. The response developed here was not to build a better model but
to build a system that keeps a model correct.

The central architectural decision was to express the system as three decoupled
Feature, Training and Inference pipelines communicating through a managed
feature store, model registry and committed published forecast documents. The
hourly workflow now couples feature refresh to forecast publication by running
inference and committing \texttt{predictions.json}; the daily workflow retrains,
evaluates and monitors the system. Because durable state lives in those services
and repository artefacts, the compute can be entirely disposable, and because the
compute is disposable it can be free. The central modelling decision was to treat the
forecast horizon as an input feature, so that one global model conditioned on
location serves 22 cities and ten lead times rather than requiring 220 artefacts.
The central operational decision was to gate promotion: a nightly retrain may
produce a new model, but it may not deploy one unless it has been measured to be
better.

\subsection{Fulfilment of the Objectives}
Each of the six objectives set in Chapter~\ref{ch:intro} was met.

\textbf{Objective 1 - Data and features.} An hourly ingestion pipeline fetches
seven pollutant and ten weather variables for 22 cities from a free, keyless
provider with retry and backoff, and a per-city feature builder expands 21 raw
columns into a 74-column engineered table with cyclical calendar encodings, wind
vector decomposition, lags at five offsets and rolling statistics over two
windows. The leakage discipline is enforced by construction - every rolling
window shifted, every family computed within a single city - and verified by
unit tests, one of which would fail if the per-city grouping were ever removed. A
resumable backfill populated the store with roughly 697,000 hourly rows spanning
January~2023 onward.

\textbf{Objective 2 - Serverless pipeline architecture.} Feature refresh and
forecast publication run hourly on free scheduled runners, while training,
evaluation and monitoring run daily. They communicate through the feature store,
model registry and committed forecast documents. A model trained on one ephemeral runner is
loaded by a different job on a different machine, which is the property on which
the whole architecture rests and which was verified directly.

\textbf{Objective 3 - Accurate multi-horizon forecasting.} A single global
gradient-boosted model with the horizon as its twenty-seventh input attains a
validation RMSE of 19.69, MAE of 12.80 and $R^2$ of 0.850, beating a persistence
baseline by 22\% overall and at every horizon from two hours to three days. Every
forecast point carries an 80\% prediction interval derived from per-horizon
empirical residual quantiles, which widens with lead time as the model's real
errors do.

\textbf{Objective 4 - MLOps discipline.} A champion-challenger gate promotes
only on a strict improvement in validation RMSE and records every decision in a
persistent leaderboard; the champion is stored in a durable versioned registry.
The gate's production record contains a first promotion, a genuine improvement
from 20.597 to 19.687 after the backfill, and a refusal of an identical model an
hour later. Accuracy is reported per horizon against a per-horizon baseline and
under a five-fold walk-forward backtest averaging 20.55, rather than as a single
averaged number.

\textbf{Objective 5 - A trust layer.} Every city's forecast carries a
plain-language explanation generated from exact Shapley attributions for its peak
hour; a what-if simulator re-predicts under user-supplied driver overrides and is
labelled explicitly as a model simulation; and the system monitors seven features
for distributional drift by Population Stability Index while separately logging
every forecast and scoring it against the air quality that subsequently arrived.

\textbf{Objective 6 - Delivery.} A typed HTTP interface serves the documented
endpoints; a five-view browser dashboard
presents the forecast, the model's promotion history, its accuracy by horizon and
its current drift status, together with global and per-city SHAP analytics and
statistical charts; and the same four grounded tools are exposed both to an
in-application advisor and over the Model Context Protocol.

\section{Contributions}
The work makes five contributions.

\textbf{A multi-city forecasting service where none existed.} Twenty-two cities
spanning every province, the capital territory, Gilgit-Baltistan and Azad Jammu
\& Kashmir receive a three-day, explained, uncertainty-quantified forecast. The
exploratory analysis reported in Section~\ref{sec:design-eda} also produced a
finding of independent interest: over the assembled record, Faisalabad's mean
index exceeds Lahore's, which is an argument against the near-exclusive attention
that Lahore receives.

\textbf{A demonstration that the full MLOps lifecycle fits inside a free tier.}
Feature store, scheduled retraining, gated promotion, durable registry, rigorous
evaluation, drift monitoring and public deployment were assembled entirely from
free and negligible-cost services. Equally importantly, the twelve concrete
failures that this constraint produced - a stalled executor freezing a blocking
insert, a statistics profiler exhausting free memory, a deployment tool honouring
an ignore file and silently omitting the model - are documented in
Table~\ref{tab:challenges} with their resolutions, and are the most transferable
part of this work.

\textbf{A worked instance of horizon-as-feature multi-horizon forecasting.} The
design collapses 220 potential artefacts into one, is shown to beat persistence
across the whole forecast window, and is shown to carry a specific and disclosed
cost at the shortest lead time.

\textbf{Honest reporting as a delivered feature.} The decaying error curve, the
backtest fold that is much worse than the others, the refused model and the
unflattering drift status all appear in the public interface rather than in an
appendix. The claim being made is not that the system is very accurate but that
the reader can see exactly how accurate it is.

\textbf{A reasoned negative result on deep sequence modelling.} A recurrent
network was implemented, evaluated on identical terms, found to be 2.3\% better at
the single horizon it was built for, and declined on grounds of coverage, training
cost and explainability - with the reasoning recorded rather than the result
suppressed.

\section{Limitations}
Six limitations bound these claims, and are stated in the same spirit as the rest
of the work.

The ground truth is a modelled atmospheric composition product rather than a
reference-grade station reading, so the system demonstrates skill at forecasting
that product. The model is trained on observed weather and deployed on forecast
weather, so the reported errors are a mild under-estimate of operational error.
The prediction intervals attain their nominal coverage by construction on the
calibration window but have not been verified prospectively. Unit testing covers
the index computation and the feature layer but leaves the gate, the intervals,
the drift computation and the API handlers to integration testing. The operational
record is short: three training runs and no matured forecast-outcome pairs yet.
And the interface's usability claims are argued rather than measured.

\section{Future Work}

\subsection{Validation Against Ground Stations}
The most valuable single extension is to obtain readings from reference-grade
monitoring stations for as many of the 22 cities as publish them, and to measure
the divergence between the modelled target and the measured one. This would either
substantiate the system's output against physical measurement or quantify a bias
that could then be corrected - a per-city calibration offset learned from the
comparison would be a natural next step.

\subsection{Prospective Interval Calibration}
The forecast log already contains everything required to measure realised interval
coverage: every issued point, its bounds, and eventually the observed value. Once
enough pairs have matured, empirical coverage should be computed per horizon and
compared with the nominal 80\%. If coverage proves to degrade during unusual
regimes, as the exchangeability assumption suggests it might, the natural remedy
is adaptive conformal prediction, in which the quantile is adjusted online in
response to recent coverage error.

\subsection{Closing the Loop Between Drift and Retraining}
At present drift is measured and retraining is scheduled, but the two are
independent: a significant drift reading does not trigger anything. A natural
extension is to make drift an input to the schedule - retraining more often, or on
a shorter and more recent window, when the index or particulate distributions shift
sharply - and to make a sustained rise in realised error raise an alert rather than
merely a table row.

\subsection{Probabilistic and Quantile Modelling}
The current interval is a post-hoc residual construction around a point forecast.
Training the model directly under a quantile loss, or fitting a small ensemble at
several quantiles, would produce intervals that adapt to the input rather than
depending only on the horizon - so that a forecast issued during a volatile
episode would be visibly less certain than one issued during a settled period.

\subsection{Pollutant-Level and Source-Aware Forecasting}
Forecasting the constituent pollutants and then composing the index, rather than
forecasting the index directly, would make the forecast more informative - a user
could be told that the problem is fine particulate matter specifically - and would
allow the gaseous sub-indices to be computed properly once the molar-mass
conversions are implemented. Incorporating exogenous signals such as satellite
fire-detection data during the crop-residue burning season would address one of
the largest sources of unmodelled variance.

\subsection{Broader Coverage and Finer Resolution}
The architecture is not specific to Pakistan or to 22 cities: adding a city is a
single line in the configuration module, and the global model requires no
retraining per location. Extending to the wider South Asian airshed, which shares
the same seasonal mechanism, is a natural direction. Finer spatial resolution
would require a different data source, since the present provider's granularity is
the ceiling on what the system can offer.

\subsection{Strengthening the Test Suite and the Interface}
The gate, the interval construction and the drift computation should be brought
under unit test, so that a regression in the system's governance is caught before
deployment rather than in production. On the interface side, a small user study
would establish whether the explanation is in fact understood, and mobile
notifications would convert the hazard alert from something a user must visit the
site to see into something that reaches them.

\section{Summary of Future Work}
The system as delivered forecasts, retrains, gates, explains and monitors itself,
and does so publicly and for free. The work that would most improve it is not more
modelling but more measurement: validating the target against physical stations,
verifying interval coverage prospectively, and letting the drift and error signals
it already computes feed back into the schedule that keeps it current. That
progression - from a system that watches itself to a system that responds to what
it sees - is the natural continuation of this thesis.
